% ============================================================================
%  TEMPLATE LaTeX — Rodapé Institucional UFMT
%  Disciplina : Algoritmos 1
%  Professora : Rafaela Souza Francisco
%  Curso      : Bacharelado em Ciência da Computação — UFMT, Cuiabá
% ============================================================================
\documentclass[12pt, a4paper]{article}

% --- Codificação e idioma ---------------------------------------------------
\usepackage[utf8]{inputenc}
\usepackage[T1]{fontenc}
% \usepackage[brazil]{babel}  % descomente se tiver texlive-lang-portuguese

% --- Fonte ------------------------------------------------------------------
% \usepackage{lmodern}        % descomente se disponível no seu sistema

% --- Layout de página -------------------------------------------------------
\usepackage[
  top=2.5cm,
  bottom=3.0cm,       % espaço extra para o rodapé
  left=2.5cm,
  right=2.5cm
]{geometry}

% --- Pacotes necessários ----------------------------------------------------
\usepackage{graphicx}          % inserir imagens
\usepackage[brazil]{babel}
\usepackage{fancyhdr}          % cabeçalho/rodapé customizado
\usepackage{xcolor}            % cores
\usepackage{hyperref}          % links clicáveis
\usepackage{tikz}              % desenhos (usado na linha decorativa)
\usepackage{amsmath, amssymb}  % símbolos matemáticos (útil para Algoritmos)
\usepackage{lastpage}          % referência à última página

\usepackage{listings}
\usepackage{xcolor}
\renewcommand{\lstlistingname}{Código}

\lstdefinestyle{cstyle}{
  language=C,
  basicstyle=\ttfamily\small,
  keywordstyle=\bfseries\color{blue!70!black},
  commentstyle=\itshape\color{green!50!black},
  stringstyle=\color{orange!80!black},
  numbers=left,
  numberstyle=\tiny\color{gray},
  numbersep=8pt,
  frame=single,
  breaklines=true,
  tabsize=4,
  showstringspaces=false,
}

\lstset{
    basicstyle=\ttfamily\small,
    keywordstyle=\color{blue}\bfseries,
    commentstyle=\color{gray}\itshape,
    stringstyle=\color{red},
    numbers=left,
    numberstyle=\tiny\color{gray},
    numbersep=8pt,
    frame=single,
    breaklines=true,
    tabsize=4,
    showstringspaces=false
}

\usetikzlibrary{shapes.geometric, arrows}
\tikzstyle{iniciofim} = [ellipse, draw, minimum width=3cm, minimum height=1cm, text centered]
\tikzstyle{processo} = [rectangle, draw, minimum width=5cm, minimum height=1cm, text centered]
\tikzstyle{seta} = [->, thick]
\tikzstyle{decisao} = [diamond, minimum width=3.5cm, minimum height=2cm, text centered, draw=black]

% --- Cores institucionais UFMT ---------------------------------------------
\definecolor{ufmtAzul}{RGB}{0, 51, 102}       % azul escuro institucional
\definecolor{ufmtVerde}{RGB}{0, 105, 62}       % verde institucional
\definecolor{ufmtCinza}{RGB}{100, 100, 100}    % cinza para texto secundário

% --- Dados do documento (EDITE AQUI) ----------------------------------------
\newcommand{\NomeAluno}{Thalles Eduardo Silva Mota}             % ← Seu nome
\newcommand{\Disciplina}{Algoritmos 1}
\newcommand{\Professora}{Profª. Rafaela Souza Francisco}
\newcommand{\Curso}{Bacharelado em Ciência da Computação}
\newcommand{\Instituicao}{UFMT — Universidade Federal de Mato Grosso}
\newcommand{\Semestre}{2026/1}
\renewcommand{\abstractname}{Resumo}
% ← Semestre atual
% Se tiver o logo em arquivo, defina o caminho abaixo:
% \newcommand{\LogoPath}{logo_ufmt.png}

% ============================================================================
%  CONFIGURAÇÃO DO RODAPÉ  (o coração do template)
% ============================================================================
\pagestyle{fancy}
\fancyhf{}                     % limpa cabeçalho e rodapé padrão

% --- Cabeçalho (opcional — descomente se quiser) ----------------------------
% \fancyhead[L]{\small\color{ufmtCinza}\Disciplina}
% \fancyhead[R]{\small\color{ufmtCinza}\NomeAluno}
% \renewcommand{\headrulewidth}{0.4pt}

% --- Rodapé -----------------------------------------------------------------
\renewcommand{\footrulewidth}{0pt}  % remove linha padrão (usaremos tikz)

\fancyfoot[C]{%
  \vspace{2pt}%
  % ── Linha decorativa com gradiente ──────────────────────────────────
  \begin{tikzpicture}[remember picture, overlay]
    \draw[
      line width=1.2pt,
      ufmtAzul
    ] ([yshift=8pt]current page.south west |- 0,0) ++(2.5cm,0)
      -- ++(\dimexpr\textwidth\relax, 0);
    % Pequeno detalhe em verde (identidade visual)
    \fill[ufmtVerde]
      ([yshift=6.5pt]current page.south west |- 0,0) ++(2.5cm,0)
      rectangle ++(4cm, 1.8pt);
  \end{tikzpicture}%
  %
  % ── Conteúdo do rodapé ──────────────────────────────────────────────
  \begin{minipage}[t]{0.12\textwidth}
    \vspace{-2pt}
    % =====================================================================
    % LOGO: Descomente a linha abaixo e comente o \fbox quando tiver o logo
    % \includegraphics[height=1.1cm]{\LogoPath}
    % =====================================================================
    % Placeholder visual enquanto não tiver o .png:
    \includegraphics[height=1.1cm]{images/UFMT - logo texto.png}%
  \end{minipage}%
  \hfill
  \begin{minipage}[t]{0.60\textwidth}
    \centering
    \vspace{-2pt}
    {\footnotesize\color{ufmtAzul}\textbf{\Instituicao}}\\[1pt]
    {\scriptsize\color{ufmtCinza}%
      \Curso\ \textbullet\ \Disciplina\ — \Semestre}\\[1pt]
    {\scriptsize\color{ufmtCinza}\Professora}
  \end{minipage}%
  \hfill
  \begin{minipage}[t]{0.15\textwidth}
    \raggedleft
    \vspace{-2pt}
    {\footnotesize\color{ufmtAzul}\textbf{\thepage}\,/\,\pageref{LastPage}}
  \end{minipage}%
}

% Estilo da primeira página (pode ser diferente)
\fancypagestyle{plain}{%
  \fancyhf{}%
  \renewcommand{\headrulewidth}{0pt}%
  \fancyfoot[C]{%
    \vspace{2pt}%
    \begin{tikzpicture}[remember picture, overlay]
      \draw[line width=1.2pt, ufmtAzul]
        ([yshift=8pt]current page.south west |- 0,0) ++(2.5cm,0)
        -- ++(\dimexpr\textwidth\relax, 0);
      \fill[ufmtVerde]
        ([yshift=6.5pt]current page.south west |- 0,0) ++(2.5cm,0)
        rectangle ++(4cm, 1.8pt);
    \end{tikzpicture}%
    \begin{minipage}[t]{0.12\textwidth}
      \vspace{-2pt}
      \fbox{\footnotesize\textbf{\color{ufmtAzul}UFMT}}%
    \end{minipage}%
    \hfill
    \begin{minipage}[t]{0.60\textwidth}
      \centering\vspace{-2pt}
      {\footnotesize\color{ufmtAzul}\textbf{\Instituicao}}\\[1pt]
      {\scriptsize\color{ufmtCinza}%
        \Curso\ \textbullet\ \Disciplina\ — \Semestre}\\[1pt]
      {\scriptsize\color{ufmtCinza}\Professora}
    \end{minipage}%
    \hfill
    \begin{minipage}[t]{0.15\textwidth}
      \raggedleft\vspace{-2pt}
      {\footnotesize\color{ufmtAzul}\textbf{\thepage}\,/\,\pageref{LastPage}}
    \end{minipage}%
  }%
}

% ============================================================================
%  DOCUMENTO DE EXEMPLO
% ============================================================================
\begin{document}

\begin{center}
  {\Large\bfseries\color{ufmtAzul} Exercícios de Aula: Expressões em C}\\[6pt]
  {\large \NomeAluno}\\[3pt]
  {\normalsize\color{ufmtCinza} \today}
\end{center}

\vspace{0.5cm}

\begin{abstract}
    \textbf{Este exercício apresenta a resolução do problema da mesada de Joãozinho 
    na linguagem C, no qual um valor inicial é acrescido de uma quantia fixa a 
    cada mês, formando uma Progressão Aritmética (PA). A implementação utiliza 
    as fórmulas do termo geral e da soma de uma PA para calcular o valor recebido 
    no último mês, o total acumulado e a média mensal. O programa demonstra 
    conceitos fundamentais de programação estruturada na linguagem C, como declaração 
    e tipagem de variáveis, leitura de dados, construção de expressões aritméticas 
    e formatação de saída, evidenciando a integração entre raciocínio matemático 
    e desenvolvimento de algoritmos.}
\end{abstract}

\noindent\textbf{Palavras-chave: Progressão Aritmética, Linguagem C, Programação 
Estruturada, Algoritmos, Expressões Aritméticas.}

\section*{Introdução}
    A aplicação de estruturas matemáticas na resolução de problemas computacionais 
    constitui um dos pilares da formação em Ciência da Computação. Traduzir um modelo 
    analítico — como uma fórmula ou uma relação de recorrência — em um algoritmo 
    funcional exige do estudante não apenas o domínio da linguagem de programação, 
    mas também a compreensão dos fundamentos matemáticos subjacentes ao problema.


Dentre os modelos matemáticos elementares com ampla aplicação prática, a Progressão 
Aritmética (PA) destaca-se pela frequência com que aparece em situações cotidianas: 
reajustes salariais fixos, planos de pagamento com parcelas crescentes e, como no 
problema aqui abordado, mesadas com acréscimos mensais constantes. Para a aplicação 
matemática do exercício proposto, será necessário o domínio da fórmula do termo geral 
de uma PA, mas também a fórmula da soma de uma PA. Ambas demonstradas no capítulo 
2 do livro citado:\cite{iezzi2013sequencias}. As fórmulas serão definidas a seguir:
\begin{equation}
  a_n = a_1 + (n-1) \cdot r
  \label{eq:termo-geral}
\end{equation}

\begin{equation}
  S_n = \frac{n \cdot (a_1 + a_n)}{2}
  \label{eq:soma-pa}
\end{equation}


Porém, a sua implementação em linguagem C mobiliza conceitos fundamentais como 
declaração e tipagem de variáveis, leitura de dados via entrada padrão, construção 
de expressões aritméticas e formatação de saída.


O presente trabalho aborda a resolução do problema da mesada de Joãozinho, no 
qual uma criança recebe um valor inicial que é acrescido de uma quantia fixa a 
cada mês. O programa desenvolvido calcula, a partir dos dados fornecidos pelo 
usuário, o valor recebido no último mês, o total acumulado ao longo do período 
e a média mensal — utilizando, respectivamente, a fórmula do termo geral, a soma 
dos termos e a média aritmética de uma PA.


O objetivo é não apenas apresentar a implementação, mas também evidenciar 
como um modelo matemático bem definido pode ser diretamente traduzido em código 
estruturado, contribuindo para a consolidação tanto dos fundamentos de algoritmos 
quanto do raciocínio matemático aplicado à programação.

\section*{Exercício 1 - Mesada do Joãozinho — Progressão Aritmética}



\begin{lstlisting}[style=cstyle, caption={Mesada do Joãozinho — Progressão 
    Aritmética}, label={lst:q2}]
/* ====================================================
   A mesada forma uma PA com:
     - Primeiro termo:  a1 = X
     - Razao:           r  = Y
     - Numero de termos: n = M
   Formulas utilizadas:
     - Ultimo mes (termo geral):  aM = X + (M-1) * Y
     - Total acumulado (soma PA): S  = M * (X + aM) / 2
     - Media:                     media = S / M
   ==================================================== */
#include <stdio.h>
int main() {
    float x, y, ultimo, total, media;
    int m;
    printf("Valor inicial da mesada (X): ");
    scanf("%f", &x);
    printf("Acrescimo mensal (Y): ");
    scanf("%f", &y);
    printf("Numero de meses (M): ");
    scanf("%d", &m);
    /* Valor no ultimo mes: a_M = X + (M-1)*Y */
    ultimo = x + (m - 1) * y;
    /* Total acumulado em M meses: S = M * (a1 + aM) / 2 */
    total = m * (x + ultimo) / 2.0;
    /* Media da mesada: media = S / M */
    media = total / m;
    printf("Ultimo mes: %.0f\n", ultimo);
    printf("Total em %d meses: %.0f\n", m, total);
    printf("Media: %.2f\n", media);
    return 0;
}
\end{lstlisting}

A solução implementada recebe do usuário três dados de entrada: o valor inicial 
da mesada $(x)$, o acréscimo mensal $(y)$ e o número de meses $(m)$. A partir 
dessas informações, o programa aplica as fórmulas da Progressão Aritmética (PA) 
para calcular três resultados.

Primeiramente, o valor recebido no último mês é obtido pela fórmula do termo 
geral da PA(Equação~\eqref{eq:termo-geral}), armazenado na variável \texttt{ultimo}. 
Em seguida, o total acumulado ao longo dos $m$ meses é calculado pela fórmula da 
soma dos termos de uma PA(Equação~\eqref{eq:soma-pa}), armazenado na variável 
\texttt{total}. Por fim, a média mensal é obtida pela divisão do total acumulado 
pelo número de meses, armazenada na variável \texttt{media}.


Essa abordagem traduz diretamente o modelo matemático em código estruturado, 
utilizando conceitos fundamentais como declaração de variáveis de tipos distintos 
(\texttt{float} para valores reais e \texttt{int} para o número de meses), leitura 
de dados via \texttt{scanf}, construção de expressões aritméticas e formatação 
de saída com \texttt{printf}. Tais conceitos são demonstrados no livro citado: 
\cite{medina2006}.

\section*{Conclusão}


O presente trabalho abordou o problema da mesada de Joãozinho, no qual uma quantia 
inicial é acrescida de um valor fixo a cada mês, caracterizando uma Progressão 
Aritmética. A solução desenvolvida na linguagem C traduziu de forma direta as 
fórmulas do termo geral e da soma de uma PA em código estruturado, demonstrando 
como um modelo matemático bem definido pode ser implementado computacionalmente 
com clareza e precisão.


A implementação mobilizou conceitos fundamentais trabalhados na disciplina de 
Algoritmos I, como a declaração de variáveis com tipagem adequada (\texttt{float} 
e \texttt{int}), a leitura de dados via \texttt{scanf}, a construção de expressões 
aritméticas compostas e a formatação de saída com \texttt{printf}. A escolha do 
tipo \texttt{float} para os valores monetários permitiu o tratamento de resultados 
não inteiros, embora em aplicações reais o uso de \texttt{double} ou de aritmética 
inteira em centavos seja recomendável para maior precisão numérica.


Além de exercitar a codificação em C, o problema evidenciou a importância da 
interdisciplinaridade entre Matemática e Computação. O reconhecimento de que 
a mesada segue o padrão de uma PA permitiu a aplicação de fórmulas fechadas — 
evitando a necessidade de estruturas de repetição — e resultando em um algoritmo 
simples, eficiente e de fácil compreensão.


Dessa forma, o exercício contribui para a consolidação tanto dos fundamentos de 
programação estruturada quanto do raciocínio matemático aplicado, reforçando a 
capacidade de identificar modelos analíticos em problemas do cotidiano e 
traduzi-los em soluções computacionais funcionais.

\bibliographystyle{abntex2-alf}  % ou plain, ieeetr, etc.
\bibliography{references}

\end{document}